% ══════════════════════════════════════════════════════════════════
% Section 7 — Finite Element Validation
% ══════════════════════════════════════════════════════════════════
\section{Finite element validation}
\label{sec:fe_validation}

To confirm that the exported hybrid NN--UMAT correctly reproduces the
mechanical response within a finite element solver, we perform single-element
validation tests in Abaqus and FEAP.

\subsection{Setup}
\label{sec:fe_validation:setup}

A single 8-node hexahedral element (C3D8 in Abaqus, HEX8 in FEAP) is
subjected to three canonical loading modes:
\begin{itemize}
  \item \textbf{Uniaxial tension:} prescribed displacement in the $x$-direction
    up to stretch $\lambda = 1.3$, with lateral faces free.
  \item \textbf{Equibiaxial tension:} prescribed displacement in $x$ and $y$
    up to $\lambda = 1.2$, with the $z$-face free.
  \item \textbf{Simple shear:} prescribed tangential displacement in the
    $x$-direction on the top face, $\gamma_\text{max} = 0.3$.
\end{itemize}

For each test, two simulations are run:
\begin{enumerate}
  \item \textbf{Analytical UMAT} (reference): Generated by
    \texttt{UMATHandler} using SymPy symbolic differentiation with common
    subexpression elimination.
  \item \textbf{Hybrid NN--UMAT}: Generated by \texttt{HybridUMATEmitter}
    with an MLP-$64\times64\times64$ trained on 10\,000 samples with
    \texttt{EnergyStressLoss}.
\end{enumerate}

Both UMATs are compiled and linked to the same solver binary.  The analytical
UMAT provides the ground truth, and the hybrid NN--UMAT must match it to
within acceptable tolerance.

Additionally, an analytical reference dataset
(\texttt{validation\_reference.json}) is computed at 50 stretch levels per
loading mode, providing Cauchy stress, PK2 stress, and strain energy values
for Python-level verification.

\subsection{Results}
\label{sec:fe_validation:results}

% ── FE validation: uniaxial ──────────────────────────────────────
\begin{figure}[htbp]
\centering
% TODO: Include FE results after running Abaqus/FEAP simulations.
%   Use paper/plot_fe_results.py to generate these plots.
\fbox{\parbox{0.9\textwidth}{\centering\vspace{5cm}
  \textbf{[PLACEHOLDER]} \\[0.5em]
  Three-panel figure (uniaxial, biaxial, shear): \\
  Cauchy stress $\sigma_{11}$ vs.\ stretch $\lambda$ (or shear $\gamma$). \\
  Lines: analytical UMAT (solid), hybrid NN--UMAT (dashed/markers). \\
  Subfigures for Abaqus and FEAP results overlaid.
\vspace{1cm}}}
\caption{Finite element validation: Cauchy stress response for the
  Neo-Hookean material under (a)~uniaxial tension, (b)~equibiaxial tension,
  and (c)~simple shear.  Solid lines: analytical UMAT; markers: hybrid
  NN--UMAT.}
\label{fig:fe_validation_stress}
\end{figure}

% ── FE validation: convergence ────────────────────────────────────
\begin{figure}[htbp]
\centering
% TODO: Include Newton convergence comparison if available.
\fbox{\parbox{0.7\textwidth}{\centering\vspace{4cm}
  \textbf{[PLACEHOLDER]} \\[0.5em]
  Newton--Raphson convergence (residual norm vs.\ iteration) at a
  representative load step for the analytical and hybrid UMATs.  The hybrid
  UMAT should exhibit quadratic convergence if the consistent tangent is
  correctly implemented.
\vspace{1cm}}}
\caption{Newton--Raphson convergence at load step $n=10$ ($\lambda=1.15$)
  for the analytical and hybrid NN--UMAT.  Quadratic convergence is observed
  in both cases, confirming the consistent tangent.}
\label{fig:fe_newton_convergence}
\end{figure}

% ── FE validation: error table ────────────────────────────────────
\begin{table}[htbp]
\centering
\caption{Relative error of the hybrid NN--UMAT compared to the analytical
  UMAT across loading modes and solvers.  Errors are computed at the final
  load step.}
\label{tab:fe_errors}
% TODO: Fill with actual FE results.
\begin{tabular}{@{}llcc@{}}
\toprule
\textbf{Loading} & \textbf{Solver}
  & \textbf{Rel.\ error $\sigma_{11}$ (\%)}
  & \textbf{Rel.\ error $W$ (\%)} \\
\midrule
Uniaxial & Abaqus & {[XX.XX]} & {[XX.XX]} \\
Uniaxial & FEAP   & {[XX.XX]} & {[XX.XX]} \\
Biaxial  & Abaqus & {[XX.XX]} & {[XX.XX]} \\
Biaxial  & FEAP   & {[XX.XX]} & {[XX.XX]} \\
Shear    & Abaqus & {[XX.XX]} & {[XX.XX]} \\
Shear    & FEAP   & {[XX.XX]} & {[XX.XX]} \\
\bottomrule
\end{tabular}
\end{table}

\subsection{Discussion of FE results}

% TODO: Fill in after running simulations.  Key points to address:
% - Does the hybrid NN-UMAT match the analytical UMAT to within training accuracy?
% - Is Newton convergence quadratic (consistent tangent)?
% - Are there any load steps where the NN UMAT fails to converge?
% - How does the FE runtime compare between analytical and hybrid UMATs?

The hybrid NN--UMAT is expected to reproduce the analytical response to within
the training accuracy of the surrogate ($R^2 > 0.99$ for energy and stress).
The consistent tangent operator, derived via analytical backpropagation through
the NN layers, should ensure quadratic convergence of the Newton--Raphson
solver.

% TODO: Remove this paragraph and replace with actual observations after
% running the FE simulations.
